% ===== PRÉAMBULE CANONIQUE — à coller VERBATIM en tête de CHAQUE .tex =====
\documentclass[11pt,a4paper]{article}
\usepackage[utf8]{inputenc}\usepackage[T1]{fontenc}\usepackage{lmodern}
\usepackage[french]{babel}
\usepackage{amsmath,amssymb,mathtools}\usepackage{icomma}
\usepackage[version=4]{mhchem}          % \ce{...} : équations & formules
\usepackage{siunitx}\sisetup{locale=FR} % \qty{}{}, \si{}, \ang{}
\usepackage{chemfig}                    % \chemfig{...} : structures développées
\usepackage{xcolor}\usepackage{colortbl}
\usepackage{tikz}\usepackage{pgfplots}\pgfplotsset{compat=1.18}
\usepackage[most]{tcolorbox}\usepackage{enumitem}\usepackage{array,booktabs}
\usepackage[a4paper,margin=2.2cm]{geometry}
\usetikzlibrary{arrows.meta,calc,angles,quotes,positioning,patterns,backgrounds,shapes.geometric}
\definecolor{primary}{RGB}{222,49,99}\definecolor{primarydark}{RGB}{150,22,55}
\definecolor{defcol}{RGB}{0,114,178}\definecolor{methcol}{RGB}{52,92,148}
\definecolor{excol}{RGB}{0,135,90}\definecolor{attncol}{RGB}{200,40,55}
\tcbset{boxbase/.style={enhanced,breakable,boxrule=0.9pt,arc=2.5pt,
  left=3mm,right=3mm,top=2.3mm,bottom=2.3mm,fonttitle=\bfseries,coltitle=white}}
% --- boîtes TUTORIEL (noms EXACTS, ne pas renommer) ---
\newtcolorbox{cle}[1][]{boxbase,colback=defcol!6,colframe=defcol,
  title={\textsc{À retenir}\ifx\relax#1\relax\relax\else\textmd{ : \itshape #1}\fi}}
\newtcolorbox{methode}[1][]{boxbase,colback=methcol!7,colframe=methcol,sharp corners,
  title={\textsc{Méthode}\ifx\relax#1\relax\relax\else\textmd{ --- \itshape #1}\fi}}
\newtcolorbox{exemplebox}[1][]{boxbase,colback=excol!5,colframe=excol,
  title={\textsc{Exemple}\ifx\relax#1\relax\relax\else\textmd{ : \itshape #1}\fi}}
\newtcolorbox{piege}[1][]{boxbase,colback=attncol!6,colframe=attncol,
  title={\textsc{Piège}\ifx\relax#1\relax\relax\else\textmd{ : \itshape #1}\fi}}
\newtcolorbox{remarque}[1][]{boxbase,colback=black!4,colframe=black!45,
  title={\textsc{Remarque}\ifx\relax#1\relax\relax\else\textmd{ : \itshape #1}\fi}}
% --- boîtes EXERCICE / CORRIGÉ (noms EXACTS, auto-comptées) ---
\newtcolorbox[auto counter]{exo}[1][]{boxbase,colback=primary!4,colframe=primary,
  title={\textsc{Exercice}~\thetcbcounter\ifx\relax#1\relax\relax\else\textmd{ --- \itshape #1}\fi}}
\newtcolorbox[auto counter]{corrige}[1][]{boxbase,colback=excol!4,colframe=excol,
  title={\textsc{Corrigé}~\thetcbcounter\ifx\relax#1\relax\relax\else\textmd{ --- \itshape #1}\fi}}
\newcommand{\R}{\mathbb{R}}\newcommand{\N}{\mathbb{N}}\newcommand{\Z}{\mathbb{Z}}\newcommand{\ds}{\displaystyle}
% (un \usepackage{fancyhdr}+en-tête est possible ; il est ignoré côté web)
% =========================================================================
\begin{document}
\sisetup{per-mode=symbol}
\begin{exo}[du nombre de CS à la précision relative]
En utilisant les ordres de grandeur typiques ($2$ CS $\sim1\%$, $3$ CS $\sim0,1\%$, $4$ CS $\sim0,01\%$),
estime l'incertitude relative d'une mesure comportant :
\begin{enumerate}[label=\alph*)]
  \item $2$ chiffres significatifs
  \item $3$ chiffres significatifs
  \item $4$ chiffres significatifs
\end{enumerate}
\end{exo}

\begin{exo}[quel instrument est le plus précis ?]
On mesure un même volume de deux façons : une éprouvette donne $(25\pm1)\ \si{\milli\litre}$, une burette
donne $(24,50\pm0,05)\ \si{\milli\litre}$.
\begin{enumerate}[label=\alph*)]
  \item Calcule l'incertitude relative de chaque mesure.
  \item Laquelle est la plus précise ? Relie ta réponse au nombre de CS.
\end{enumerate}
\end{exo}

\begin{exo}[une balance déréglée]
Une balance biaisée pèse cinq fois un objet de masse \emph{vraie} $\qty{48,00}{\gram}$ et affiche :
$50,01\,;\ 49,99\,;\ 50,00\,;\ 50,02\,;\ 50,00\ \si{\gram}$.
\begin{enumerate}[label=\alph*)]
  \item La balance est-elle \textbf{fidèle} ?
  \item Est-elle \textbf{juste} ?
  \item Que faut-il faire pour corriger le problème ?
\end{enumerate}
\end{exo}

\begin{exo}[une écriture incohérente]
Un étudiant note le résultat $m=(5,2\pm0,01)\ \si{\gram}$.
\begin{enumerate}[label=\alph*)]
  \item Qu'y a-t-il d'incohérent entre la valeur et son incertitude ?
  \item Propose deux écritures cohérentes selon la précision réelle de la balance.
\end{enumerate}
\end{exo}

\begin{exo}[de la précision relative au nombre de CS]
Réciproquement, à combien de chiffres significatifs peut-on \emph{légitimement} afficher un résultat
connu avec une incertitude relative de :
\begin{enumerate}[label=\alph*)]
  \item $\sim 1\%$
  \item $\sim 0,01\%$
\end{enumerate}
\end{exo}

\end{document}
